\section*{第 \ref{ch:g11:subduction-zones} 章——俯冲带}

\begin{solution}{exo:g11:subduction-zones:1}
\omterm{def:g11:subduction-zones:subduction}{俯冲}是指一块\omterm{def:g11:ocean-ridges:oceanic-lithosphere}{大洋岩石圈}板块
插入另一板块之下、重新没入地幔的过程；\omterm{def:g11:subduction-zones:subduction}{俯冲带}是发生这一过程的汇聚型板块边界，
以海沟为标志。\omterm{def:g11:ocean-ridges:oceanic-lithosphere}{大洋岩石圈}会下沉：一旦变老，
它便变冷、变厚，密度大于其下方的\omterm{def:g11:lithosphere-asthenosphere:lithosphere}{软流圈}。
陆壳太轻，不会下沉。
\end{solution}

\begin{solution}{exo:g11:subduction-zones:2}
（1）例如汤加–斐济地区、从日本到马里亚纳群岛的西太平洋、
印度尼西亚和南美洲。（2）总是位于上覆板块一侧，
从不位于海洋一侧。（3）浅源：图中大西洋中脊的地震，
均浅于 \qty{70}{km}。
\end{solution}

\begin{solution}{exo:g11:subduction-zones:3}
$7031 \times 0.21 \approx 1480$ 个深于 \qty{70}{km} 的地震；
$7031 \times 0.069 \approx 485$ 个深于 \qty{300}{km} 的地震。
\end{solution}

\begin{solution}{exo:g11:subduction-zones:4}
\omterm{def:g11:subduction-zones:eclogite}{蓝片岩}是由大洋地壳的玄武岩
在相对低温下埋至高压条件下形成的变质岩，因含蓝闪石而呈蓝色；
\omterm{def:g11:subduction-zones:eclogite}{榴辉岩}则是玄武岩和\omterm{def:g11:ocean-ridges:gabbro}{辉长岩}在地幔深处形成的致密变质岩，
由红色石榴子石和绿色辉石组成。两者都源自大洋地壳，
并保持其化学组成；二者的差别在于
在新的压力和温度下以固态形成的矿物，
以及密度。
\end{solution}

\begin{solution}{exo:g11:subduction-zones:5}
洋底：老的\omterm{def:g11:lithosphere-asthenosphere:lithosphere}{岩石圈}向海沟运动。海沟：板块弯曲，
开始下沉。弧前：从下沉板块上刮落的沉积物堆积；其下，
两板块沿板块边界相互错动，
大地震即发生于此。火山弧：位于\omterm{def:g11:subduction-zones:slab}{板片}上方的一排火山。
弧后：在某些火山弧的后方，上覆板块被拉伸。
\end{solution}

\begin{solution}{exo:g11:subduction-zones:6}
$\tan\delta = (508 - 135)/(800 - 300) = 373/500 \approx 0.746$，故
$\delta \approx \ang{37}$。
\end{solution}

\begin{solution}{exo:g11:subduction-zones:7}
（1）$\tan\ang{45} = 1$：深 \qty{300}{km}。（2）距海沟
\qty{670}{km} 处。（3）\omterm{def:g11:subduction-zones:slab}{板片}长度 $670\sqrt{2} \approx \qty{948}{km}$。
\end{solution}

\begin{solution}{exo:g11:subduction-zones:8}
（1）$(74.6 + 70.0 + 65.5 + 63.1 + 52.6)/5 \approx \qty{65.2}{mm}$ 每年；
$(47.6 + 31.0 + 29.2 + 22.8 + 19.8 + 15.2 + 14.9 + 10.6)/8 \approx
\qty{23.9}{mm}$ 每年；比值 $65.2/23.9 \approx 2.7$。（2）板块
在有\omterm{def:g11:subduction-zones:slab}{板片}下垂时运动最快：下沉\omterm{def:g11:subduction-zones:slab}{板片}的拉力是
驱动板块运动的最强作用力，正如福赛思与宇田所得出的结论。
\end{solution}

\begin{solution}{exo:g11:subduction-zones:9}
（1）\qty{83}{mm} 即 \qty{0.083}{m}：$50/0.083 \approx 600$ 年和
$60/0.083 \approx 720$ 年的板块运动。（2）从 869 到 2011 共 1142
年，与之同一数量级。若板块边界在两次大地震之间保持锁定，
则 2011 年的错动释放了数百年积累的运动量。
\end{solution}

\begin{solution}{exo:g11:subduction-zones:10}
（1）日本：$6400 - 3200\,e^{-133/62.8} \approx \qty{6015}{m}$；汤加：
$6400 - 3200\,e^{-99/62.8} \approx \qty{5740}{m}$；安第斯：
$2500 + 350\sqrt{49} = \qty{4950}{m}$。（2）到达日本的海底是
最古老的，因而\omterm{def:g11:lithosphere-asthenosphere:lithosphere}{岩石圈}最冷、密度最大，海底也最深。
\end{solution}

\begin{solution}{exo:g11:subduction-zones:11}
（1）$3380 \times 1.06 \approx \qty{3583}{kg/m^3}$。（2）
$(3583 - 2900)/2900 \approx 0.24$：约升高 \qty{24}{\percent}。
（3）变化前，地壳比地幔轻
$3380 - 2900 = \qty{480}{kg/m^3}$；变化后则约大
\qty{203}{kg/m^3}。它不再托住\omterm{def:g11:subduction-zones:slab}{板片}，反而将其拉下。
\end{solution}

\begin{solution}{exo:g11:subduction-zones:12}
（1）$t = 6 \times 10^6 \times 3.16 \times 10^7 \approx 1.9 \times 10^{14}$~秒；
$\sqrt{10^{-6} \times 1.9 \times 10^{14}} \approx \qty{1.4e4}{m}$，约
\qty{14}{km}。（2）远小于 \qty{100}{km}：在 600 万年内，只有
\omterm{def:g11:subduction-zones:slab}{板片}外壳升温，其内部仍保持低温。（3）冷岩石
坚硬且能破裂，而周围灼热的地幔则流动；深源
地震发生在冷的\omterm{def:g11:subduction-zones:slab}{板片}之中。
\end{solution}

\begin{solution}{exo:g11:subduction-zones:13}
（1）$\sqrt{1017^2 + 670^2} \approx \qty{1218}{km}$。（2）\qty{64}{mm} 每年
即 \qty{64}{km} 每百万年：$1218/64 \approx 19$ 百万年。
（3）比汤加下方约长三倍：安第斯下方的\omterm{def:g11:subduction-zones:slab}{板片}下沉更慢，
倾角也更缓。
\end{solution}

\begin{solution}{exo:g11:subduction-zones:14}
（1）沿海沟取宽 \qty{1}{m} 的一片，其体积为
$10^5 \times 10^6 \times 1 = \qty{1e11}{m^3}$；其过剩质量为
$36 \times 10^{11} = \qty{3.6e12}{kg}$。（2）其过剩重量为
$3.6 \times 10^{12} \times 9.8 \approx \qty{3.5e13}{N}$，即海沟
每米的过剩重量。（3）\omterm{def:g11:subduction-zones:slab}{板片}与板块的其余部分相连，板块为刚性，
这一作用力传递至整个板块，使之被拖向
海沟，况且海沟延伸数千千米。
\end{solution}

\begin{solution}{exo:g11:subduction-zones:15}
（1）$(23.6 - 16.5) \times 111 \approx \qty{790}{km}$。（2）速率增加
$237 - 104 = \qty{133}{mm}$ 每年，分布在 \qty{790}{km} 上，约每百千米
\qty{17}{mm} 每年。（3）汇聚速率即
太平洋板块的速度减去火山弧的速度。由于太平洋板块的速度
在两点处大致相同，火山弧本身必定朝太平洋板块方向运动，
且越往北越快：其后方的弧后正在
被拉伸，北部比南部更甚。
\end{solution}

\begin{solution}{pb:g11:subduction-zones:1}
\textbf{1.} 约 \qty{671}{km}。

\textbf{2.} 在海洋一侧，板块尚未下沉：此处的\omterm{def:g11:lithosphere-asthenosphere:lithosphere}{岩石圈}
靠近地表。深源地震发生在\omterm{def:g11:subduction-zones:slab}{板片}内部，而\omterm{def:g11:subduction-zones:slab}{板片}
仅位于上覆板块之下。

\textbf{3.} $z(300) = -24 + 273 = \qty{249}{km}$ 和
$z(700) = -24 + 637 = \qty{613}{km}$；$\tan\delta = 364/400 = 0.91$，故
$\delta \approx \ang{42}$。

\textbf{4.} $x = (100 + 24)/0.91 \approx \qty{136}{km}$ 和
$x = (400 + 24)/0.91 \approx \qty{466}{km}$。

\textbf{5.} 若从海沟处出发、倾角如此的直线，将
经过 $z = 0$ 对应的 $x = 0$；\omterm{def:g11:subduction-zones:slab}{板片}位于其上方。\omterm{def:g11:subduction-zones:slab}{板片}并非直线：靠近
海沟处倾角较缓，向深处则逐渐变陡。

\textbf{6.} 每年 \qty{16.8}{cm}，即每百万年 \qty{168}{km}。

\textbf{7.} 快 $168/69 \approx 2.4$ 倍。

\textbf{8.} 跨海沟的汇聚速率等于太平洋板块的速度
减去另一侧火山弧的速度。若火山弧朝太平洋板块方向运动，
正如其后方弧后被拉伸时所发生的那样，两者速度相加，
汇聚速率便大于太平洋板块
自身的速度。

\textbf{9.} $6400 - 3200\,e^{-99/62.8} \approx \qty{5740}{m}$。

\textbf{10.} 约 6000 万年后，\omterm{def:g11:ocean-ridges:oceanic-lithosphere}{大洋岩石圈}整体上
比其下方的\omterm{def:g11:lithosphere-asthenosphere:lithosphere}{软流圈}更致密；到 9900 万年时更甚，
边缘一旦下沉，其重量便会将其余部分拽入。

\textbf{11.} $(670 + 24)/0.91 \approx \qty{763}{km}$。

\textbf{12.} $\sqrt{763^2 + 670^2} \approx \qty{1015}{km}$。

\textbf{13.} $1015/168 \approx 6.0$ 百万年。

\textbf{14.} $1015/69 \approx 14.7$ 百万年。

\textbf{15.} $3380 \times 1.06 \approx \qty{3583}{kg/m^3}$：地壳原本比地幔轻，
如今变得比地幔更致密，整个\omterm{def:g11:subduction-zones:slab}{板片}的拉力随之
增大。

\textbf{16.} $t \approx 6.0 \times 10^6 \times 3.16 \times 10^7 \approx
1.9 \times 10^{14}$~s，且 $\sqrt{\kappa t} \approx \qty{14}{km}$。

\textbf{17.} \qty{14}{km} 远小于 \qty{100}{km}：只有\omterm{def:g11:subduction-zones:slab}{板片}外层
已被加热，其内部在抵达
\qty{670}{km} 时仍然寒冷。

\textbf{18.} 冷的\omterm{def:g11:subduction-zones:slab}{板片}仍足够坚硬，直至其尖端都能断裂，
而其周围的热地幔则在流动：汤加之下的地震
发生在\omterm{def:g11:subduction-zones:slab}{板片}内部，深度可达约 \qty{671}{km}。

\textbf{19.} 地震波显示，地幔在约
\qty{670}{km} 处发生急剧变化：波速在此处跃变。

\textbf{20.} 从海沟到
\qty{670}{km} 约需 \textbf{600 万年}。
\end{solution}
